% !TeX root = ../main.tex
% Add the above to each chapter to make compiling the PDF easier in some editors.

\chapter{Evaluation}\label{chapter:evaluation}

This chapter describes the evaluation methodology, metrics, datasets, and results for the proposed multi-frame calibration and pose estimation system, comparing trained models against the AnyCam and AnyCalib baselines.

\section{Evaluation Methodology}\label{sec:eval-methodology}

\subsection{Evaluation Principle}

All training in this work is \textbf{fully self-supervised}: no ground truth labels for camera poses or intrinsics are used during training. Ground truth is used exclusively for evaluation and hyperparameter validation. This matches AnyCam's original training paradigm and enables scaling to large, unlabeled video datasets where ground truth is unavailable or expensive to obtain.

The self-supervised training relies on geometric constraints provided by flow reprojection: predicted camera motion and depth should induce optical flow that matches observed flow from a frozen UniMatch model. Ground truth serves only as an external validation signal to measure how well the learned geometry aligns with the true camera parameters.

\subsection{What is Evaluated}

The evaluation focuses on three primary aspects:

\begin{enumerate}
    \item \textbf{Pose accuracy}: How well does the model predict relative camera motion (rotation and translation) between frames?
    \item \textbf{Calibration accuracy and consistency}: How well does the model predict camera intrinsic parameters, and how consistent are predictions across frames?
    \item \textbf{Comparison baselines}: How does the model compare to AnyCam (original model) and AnyCalib (standalone per-frame averaging)?
\end{enumerate}

\textbf{Evaluation goal}: The primary objective is to demonstrate that the proposed method achieves \emph{more accurate and consistent calibration than AnyCalib} while simultaneously \emph{improving pose accuracy beyond AnyCam}. Raw model comparisons are performed without post-optimisation (e.g., bundle adjustment), as any post-processing applied to AnyCam could equally be applied to our method.

For each aspect, multiple metrics are computed to provide a comprehensive view of performance across different error types and statistical summaries (mean, median, variance).

\section{Evaluation Metrics}\label{sec:eval-metrics}

\subsection{Pose Metrics}\label{sec:pose-metrics}

Three complementary metrics are used to evaluate pose prediction accuracy:

\subsubsection{Rotation Error (SO(3) Geodesic Distance)}

Rotation error is computed as the geodesic distance on SO(3) between predicted and ground truth rotation matrices. Given predicted rotation $\mathbf{R}_{\text{pred}}$ and ground truth rotation $\mathbf{R}_{\text{gt}}$, the error in degrees is:
\begin{equation}
    \theta_{\text{rot}} = \arccos\left(\frac{\text{trace}(\mathbf{R}_{\text{pred}}^\top \mathbf{R}_{\text{gt}}) - 1}{2}\right) \cdot \frac{180}{\pi}
\end{equation}

This metric measures the angle of the residual rotation needed to align the prediction with ground truth. Both mean and median errors are reported.

\subsubsection{Translation Direction Error}

Translation direction error is computed as the angular difference between predicted and ground truth translation vectors:
\begin{equation}
    \theta_{\text{trans}} = \arccos\left(\frac{\mathbf{t}_{\text{pred}} \cdot \mathbf{t}_{\text{gt}}}{\|\mathbf{t}_{\text{pred}}\| \|\mathbf{t}_{\text{gt}}\|}\right) \cdot \frac{180}{\pi}
\end{equation}

This metric is scale-invariant, appropriate for monocular systems where absolute scale cannot be recovered.

\subsubsection{Average Trajectory Error (ATE)}

Average trajectory error quantifies the overall drift in camera path prediction. Using the evo library~\parencite{gruber2017evo}, we:
\begin{enumerate}
    \item Align predicted and ground truth trajectories using Procrustes alignment (SE(3) + scale)
    \item Compute the Euclidean distance between aligned camera positions at each frame
    \item Report the mean, median, and RMSE of these distances
\end{enumerate}

ATE provides a global measure of trajectory quality, complementing the per-frame rotation and translation metrics. Lower ATE indicates better consistency across longer sequences.

\subsection{Calibration Metrics}\label{sec:calib-metrics}

Calibration is evaluated using both accuracy (when GT available) and consistency (always measurable):

\subsubsection{Mean Absolute Error (MAE) and MAPE}

When ground truth calibration is available, MAE measures absolute error in pixels:
\begin{equation}
    \text{MAE}(p) = \frac{1}{N} \sum_{i=1}^N |p_{\text{pred}}^{(i)} - p_{\text{gt}}^{(i)}| \quad \text{where } p \in \{f_x, f_y, c_x, c_y\}
\end{equation}

MAPE measures relative error as a percentage, enabling fair comparison across image resolutions:
\begin{equation}
    \text{MAPE}(p) = \frac{100}{N} \sum_{i=1}^N \frac{|p_{\text{pred}}^{(i)} - p_{\text{gt}}^{(i)}|}{|p_{\text{gt}}^{(i)}|}
\end{equation}

\subsubsection{Focal Length Consistency (Pseudo-Benchmark for Datasets Without GT)}

For datasets without ground truth calibration, we measure \textbf{prediction consistency} across frames as a proxy for quality. This assumes:
\begin{itemize}
    \item Camera intrinsics are constant across a video sequence (no zooming or autofocus)
    \item More consistent predictions indicate better calibration quality
\end{itemize}

The consistency metric is the coefficient of variation (normalized standard deviation):
\begin{equation}
    \text{CV}(p) = \frac{\sigma(p)}{\mu(p)} \times 100\% \quad \text{where } \sigma = \text{std dev}, \mu = \text{mean}
\end{equation}

Lower CV indicates more stable predictions. We compare:
\begin{enumerate}
    \item \textbf{AnyCalib per-frame}: CV of raw per-frame predictions (baseline)
    \item \textbf{Our method}: CV of MCT aggregated predictions (should be lower)
\end{enumerate}

This pseudo-benchmark enables evaluation on datasets where GT calibration is unavailable.

\section{Datasets}\label{sec:datasets}

Training datasets are described in Section~\ref{sec:training-data}. This section covers the datasets used for evaluation, which provide ground truth labels for measuring model performance.

\subsection{Evaluation Datasets}\label{sec:eval-datasets}

Three datasets are used for final evaluation, covering diverse indoor and outdoor scenarios:

\subsubsection{MPI Sintel}

MPI Sintel~\parencite{butler2012sintel} is a synthetic dataset rendered from an open-source animated film, providing ground truth camera poses, intrinsics, depth maps, and optical flow. The dataset features diverse scene types including outdoor markets, narrow alleys, and temple interiors with challenging dynamic content (moving characters, cloth, particles). Because all ground truth is derived from the rendering pipeline, Sintel supports \textbf{both pose and calibration evaluation} with exact annotations, making it ideal for comprehensive assessment.

\subsubsection{TUM-RGBD}

The TUM RGB-D benchmark~\parencite{sturm2012tumrgbd} contains indoor sequences captured with handheld RGB-D cameras, with high-precision ground truth trajectories from an external motion capture system and known camera intrinsics. The dataset includes sequences with varying motion dynamics---from slow desktop movements to fast handheld walking---providing a rigorous test of pose estimation robustness. Used for \textbf{pose and calibration evaluation}.

\subsubsection{KITTI}

The KITTI Vision Benchmark~\parencite{geiger2012kitti} provides automotive driving sequences captured with calibrated stereo cameras and a GPS/IMU system for ground truth poses. Sequences feature forward-facing camera motion through urban and highway environments with dynamic traffic. Used for \textbf{pose evaluation}; calibration evaluation on KITTI is discussed as a limitation in Section~\ref{sec:kitti-limitation}.

\section{Results}\label{sec:fullscale-results}

This section presents the evaluation results across all training phases, comparing the proposed method against the AnyCam and AnyCalib baselines.

\subsection{Training Convergence}\label{sec:training-convergence}

Training proceeded through all three phases on the combined RealEstate10K, YouTube VOS, WalkingTours, and OpenDV datasets (approximately 82,000 frames across 1,774 videos, preprocessed into 77,000 training sequences).

\textbf{Phase A} (pose head initialisation) trained only the pose head using static AnyCalib calibration with the flow reprojection objective. Training used a learning rate of $10^{-4}$ with the AdamW optimiser, batch size 20, and 4-frame windows ($\text{max\_ahead}=3$). Only the pose head was trainable; the DINOv2-small backbone and all other components remained frozen. Figure~\ref{fig:training-curves-a} shows stable convergence over 5 epochs, with validation loss closely tracking training loss.

\begin{figure}[htpb]
    \centering
    \includegraphics[width=0.7\textwidth]{figures/training_loss_phase_a.png}
    \caption[Phase A training convergence]{Phase~A training and validation loss (flow reprojection). The pose head converges rapidly within the first two epochs, with validation loss tracking training loss closely, indicating no overfitting.}
    \label{fig:training-curves-a}
\end{figure}

\textbf{Phase B} (MCT pre-training) trained only the MCT module (${\sim}25$M parameters) using the calibration reprojection loss against averaged AnyCalib pseudo ground truth. All backbone components remained frozen during this phase. Figure~\ref{fig:training-curves-b} shows that training loss decreases steadily over 10 epochs from 19.6 to 2.4. Validation loss is noisy due to the limited validation set size, but the training loss trend confirms that MCT successfully learns to aggregate spatial features.

\begin{figure}[htpb]
    \centering
    \includegraphics[width=0.7\textwidth]{figures/training_loss_phase_b.png}
    \caption[Phase B training convergence]{Phase~B training and validation loss (calibration reprojection against AnyCalib pseudo ground truth). Training loss decreases steadily; validation loss fluctuates due to limited validation set size but does not diverge.}
    \label{fig:training-curves-b}
\end{figure}

\textbf{Phase C} (end-to-end training) jointly optimised the pose and calibration components. The trainable parameters comprised: the pose neck (reassemble stage, fusion stage, interframe attention, sequence token; ${\sim}2.5$M parameters), the pose head (${\sim}21$K parameters), and the MCT (${\sim}25$M parameters)---approximately 27.5M trainable parameters in total. Both DINOv2 backbones (small and ViT-L/14), the DPT decoder, UniDepth, and UniMatch remained frozen throughout. Training used a learning rate of $10^{-4}$, batch size 20, and 5 epochs on an NVIDIA H100 GPU. The combined loss included flow reprojection and a calibration anchor ($\lambda_{\text{calib}} = 10^{-4}$).

Figure~\ref{fig:training-curves-c} shows the Phase~C loss components. The flow reprojection loss decreases steadily, confirming that end-to-end training improves pose quality. The calibration anchor loss increases moderately, which is expected: as the model optimises flow reprojection, it learns to adjust calibration predictions away from the AnyCalib pseudo ground truth where doing so improves the self-supervised objective. Validation loss improves consistently, indicating generalisation.

Two variants of Phase~C were explored: one that additionally unfroze the DINOv2 ViT-L/14 backbone during training, and one that kept both backbones frozen. The unfrozen-backbone variant diverged during training, with losses becoming unstable as the large backbone's pretrained representations were disrupted by the comparatively weak calibration training signal. The frozen-backbone variant produced stable convergence and was used for all reported evaluations. A further Phase~D experiment attempted pose-only fine-tuning from the Phase~C checkpoint (freezing the calibration transformer and training only the pose head), but this degraded performance, confirming that the joint training of Phase~C already achieves a good balance between pose and calibration optimisation.

\begin{figure}[htpb]
    \centering
    \includegraphics[width=\textwidth]{figures/training_loss_phase_c.png}
    \caption[Phase C training convergence]{Phase~C training loss components over 5 epochs. Left: total loss and validation loss both decrease, confirming stable end-to-end training. Centre: flow reprojection loss decreases steadily. Right: calibration anchor loss increases as the model learns to deviate from AnyCalib pseudo ground truth where doing so improves flow reprojection---an expected and desirable behaviour.}
    \label{fig:training-curves-c}
\end{figure}

\subsection{Pose Estimation Results}\label{sec:pose-results}

Table~\ref{tab:pose-results} presents the pose estimation results across all evaluation datasets, comparing our method against the AnyCam baseline. All comparisons are on raw model outputs without post-optimisation (e.g., bundle adjustment), ensuring a fair comparison of the learned representations.

\begin{table}[htpb]
    \centering
    \caption[Pose estimation results]{Pose estimation results across evaluation datasets ($N=200$ sequences, 600 pose pairs per dataset). Rotation error (SO(3) geodesic distance in degrees) and translation direction error (degrees). Best results in bold. Improvement percentages shown for reference.}
    \label{tab:pose-results}
    \footnotesize
    \renewcommand{\arraystretch}{1.25}
    \begin{tabular}{@{}ll cc @{}}
        \toprule
        Dataset & Method & Rot.\ Error ($^\circ$) $\downarrow$ & Trans.\ Dir.\ ($^\circ$) $\downarrow$ \\
        & & mean / median & mean / median \\
        \midrule
        \multirow{3}{*}{Sintel} & AnyCam & 0.67 / 0.21 & 89.30 / 86.53 \\
        & Ours & \textbf{0.60 / 0.19} & \textbf{64.95 / 52.33} \\
        & & {\scriptsize $-$10.1\% / $-$8.2\%} & {\scriptsize $-$27.3\% / $-$39.5\%} \\
        \midrule
        \multirow{3}{*}{TUM-RGBD} & AnyCam & 2.03 / 1.23 & 93.04 / 97.05 \\
        & Ours & \textbf{1.39 / 0.71} & \textbf{71.94 / 66.90} \\
        & & {\scriptsize $-$31.7\% / $-$42.6\%} & {\scriptsize $-$22.7\% / $-$31.1\%} \\
        \midrule
        \multirow{3}{*}{KITTI} & AnyCam & 0.56 / 0.26 & 91.14 / 91.45 \\
        & Ours & \textbf{0.54 / 0.25} & \textbf{77.20 / 70.93} \\
        & & {\scriptsize $-$2.0\% / $-$2.4\%} & {\scriptsize $-$15.3\% / $-$22.4\%} \\
        \bottomrule
    \end{tabular}
\end{table}

Our method achieves improved pose accuracy compared to AnyCam across all three evaluation datasets. The improvements are strongest on \textbf{TUM-RGBD}, where rotation error decreases by 31.7\% (mean) and 42.6\% (median), and translation direction error improves by 22.7\% (mean) and 31.1\% (median). On \textbf{Sintel}, rotation improves by 10.1\% and translation direction by 27.3--39.5\%. On \textbf{KITTI}, rotation improvements are modest (2.0\%) but translation direction still improves by 15.3--22.4\%.

The per-dataset pattern is consistent: translation direction error benefits most from multi-frame consistency enforcement, as composed flow losses provide additional geometric constraints that reduce translational drift. Rotation improvements are more modest on datasets where AnyCam's pairwise rotation predictions are already strong (Sintel, KITTI), but substantial on TUM-RGBD where the handheld camera motion with frequent direction changes benefits from multi-frame context. The consistent improvements across synthetic (Sintel), indoor real (TUM-RGBD), and automotive (KITTI) domains confirm that the staged training pipeline generalises across scene types.

Figures~\ref{fig:rotation-error-hist} and~\ref{fig:translation-error-hist} present per-pair error distributions on a representative Sintel sequence, showing how our method concentrates errors at lower values compared to the AnyCam baseline.

\begin{figure}[htpb]
    \centering
    \includegraphics[width=0.85\textwidth]{figures/rotation_error_histogram.png}
    \caption[Rotation error distribution]{Per-pair rotation error distribution on Sintel \texttt{alley\_1} (147 pose pairs). Our method (median $0.08^\circ$) concentrates errors at lower values compared to AnyCam (median $0.11^\circ$). Both methods achieve sub-degree accuracy on this sequence, but our method produces a tighter distribution.}
    \label{fig:rotation-error-hist}
\end{figure}

\begin{figure}[htpb]
    \centering
    \includegraphics[width=0.85\textwidth]{figures/translation_error_histogram.png}
    \caption[Translation direction error distribution]{Per-pair translation direction error distribution on Sintel \texttt{alley\_1}. Our method (median $34.4^\circ$) substantially outperforms AnyCam (median $107.2^\circ$), demonstrating that improved calibration directly benefits translation direction estimation.}
    \label{fig:translation-error-hist}
\end{figure}

\subsection{Calibration Accuracy and Consistency Results}\label{sec:calib-results}

Calibration is evaluated using two complementary approaches: accuracy against ground truth (when available) and prediction consistency across frames.

\subsubsection{Calibration Accuracy (Datasets with GT)}

Table~\ref{tab:calib-accuracy} compares calibration accuracy on datasets with ground truth intrinsics, measuring MAE and MAPE for each intrinsic parameter.

\begin{table}[htpb]
    \centering
    \caption[Calibration accuracy results]{Focal length accuracy on datasets with ground truth intrinsics ($N=200$ sequences per dataset). $f_x$ MAE in pixels, combined $f$ MAPE as percentage. Three methods compared: AnyCam's 32-candidate selection, AnyCalib per-frame averaging, and our MCT-aggregated predictions. Best results in bold.}
    \label{tab:calib-accuracy}
    \footnotesize
    \renewcommand{\arraystretch}{1.25}
    \begin{tabular}{@{}ll cc @{}}
        \toprule
        Dataset & Method & $f_x$ MAE (px) $\downarrow$ & $f$ MAPE (\%) $\downarrow$ \\
        \midrule
        \multirow{3}{*}{Sintel} & AnyCam (32-cand.) & 502.5 & 68.2 \\
        & AnyCalib (per-frame) & 329.9 & 30.7 \\
        & Ours (MCT) & \textbf{300.3} & \textbf{27.8} \\
        \midrule
        \multirow{3}{*}{TUM-RGBD} & AnyCam (32-cand.) & 234.1 & 63.7 \\
        & AnyCalib (per-frame) & 43.0 & 11.7 \\
        & Ours (MCT) & \textbf{30.3} & \textbf{7.9} \\
        \bottomrule
    \end{tabular}
\end{table}

MCT-aggregated calibration achieves lower MAPE than both AnyCalib per-frame averaging and AnyCam's 32-candidate system on all evaluated datasets. On \textbf{TUM-RGBD}, MCT reduces focal length MAPE by 32.5\% relative to AnyCalib (from 11.7\% to 7.9\%) and by 87.6\% relative to AnyCam. On \textbf{Sintel}, the improvement over AnyCalib is 9.5\% (from 30.7\% to 27.8\%). The absolute MAPE values on Sintel are higher than on TUM-RGBD, reflecting the greater diversity of camera configurations in the synthetic dataset.

The improvement over AnyCalib confirms the core hypothesis that feature-level aggregation preserves geometric information lost in scalar averaging. By combining spatial features before decoding, MCT can identify and weight the most reliable geometric cues across frames, rather than treating all predictions equally. The large improvement over AnyCam's 32-candidate system further validates the architectural choice to replace discrete candidate search with direct prediction.

\subsubsection{Calibration Stability}

Beyond accuracy, temporal stability of calibration predictions is important for downstream applications. AnyCalib processes each frame independently, producing per-frame estimates that can vary significantly even within a single video where the true intrinsics are constant. MCT aggregates features across frames before decoding, which should inherently produce more stable predictions.

Figure~\ref{fig:calib-stability} illustrates this on a representative Sintel sequence (\texttt{alley\_1}, 49 windows). The AnyCalib per-frame predictions (individual points) scatter around the ground truth with noticeable frame-to-frame variation, while AnyCam's 32-candidate system produces occasional large outliers. Our MCT predictions (connected line) are substantially smoother, tracking closer to the ground truth focal length. This temporal stability is a direct consequence of the multi-frame feature aggregation: by attending across frames at the feature level, MCT produces predictions that reflect the consensus of multiple observations rather than relying on a single frame's features.

\begin{figure}[htpb]
    \centering
    \includegraphics[width=0.95\textwidth]{figures/focal_length_over_time.png}
    \caption[Calibration stability across frames]{Focal length predictions over consecutive windows on Sintel \texttt{alley\_1} (GT $f_x = 530.1$\,px). MCT predictions (solid line) remain close to the ground truth with smooth temporal variation, while AnyCalib per-frame predictions (dots) scatter with higher variance. AnyCam's 32-candidate selection produces occasional large outliers. Mean absolute errors: MCT $28.5$\,px, AnyCalib $40.4$\,px, AnyCam $362.1$\,px.}
    \label{fig:calib-stability}
\end{figure}

Figure~\ref{fig:focal-error-hist} further illustrates the calibration accuracy distribution, showing how MCT concentrates its predictions near the ground truth focal length while AnyCam's 32-candidate system spreads errors across a wide range.

\begin{figure}[htpb]
    \centering
    \includegraphics[width=0.85\textwidth]{figures/focal_length_error_histogram.png}
    \caption[Focal length error distribution]{Distribution of focal length errors on Sintel \texttt{alley\_1}. MCT errors (mean $28.5$\,px) are concentrated near zero, AnyCalib errors (mean $40.4$\,px) show moderate spread, while AnyCam's 32-candidate system (mean $362.1$\,px) produces a wide distribution with large outliers.}
    \label{fig:focal-error-hist}
\end{figure}

\subsection{Sequence Length Selection}\label{sec:multiframe-generalization}

The choice of sequence length ($\text{max\_ahead}=3$, corresponding to 4-frame windows) was informed by the preliminary hyperparameter sweep described in Appendix~\ref{chapter:preliminary-results}. Table~\ref{tab:max-ahead-sweep} shows that $\text{max\_ahead}=3$ achieves the best translation direction error while maintaining competitive rotation accuracy. Shorter sequences ($k=2$) underutilise multi-frame information, while longer sequences ($k \geq 4$) suffer from compounding flow composition errors.

All full-scale training was conducted with $\text{max\_ahead}=3$ based on these preliminary findings. The consistent improvements observed across all three evaluation datasets (Table~\ref{tab:pose-results}) confirm that this sequence length generalises well from the preliminary validation setting to large-scale training and diverse evaluation domains.

\subsection{Ablation Studies}\label{sec:ablations}

Due to the computational cost of full-scale training (each Phase~C run requires multiple days on an H100 GPU), systematic ablations over all hyperparameters were not feasible at full scale. Instead, the key design choices---sequence length and calibration loss weight---were validated during preliminary experiments on the Objectron dataset (Appendix~\ref{chapter:preliminary-results}), and the selected values were used for all full-scale training runs.

\subsubsection{Sequence Length (max\_ahead)}

The preliminary sweep over $\text{max\_ahead} \in \{2, 3, 4, 5, 6, 7\}$ (Table~\ref{tab:max-ahead-sweep}) identified $\text{max\_ahead}=3$ as optimal, achieving the best translation direction error ($64.13^\circ$ mean) while maintaining competitive rotation accuracy ($1.37^\circ$ mean). The consistent improvements observed across all three full-scale evaluation datasets (Table~\ref{tab:pose-results}) validate this choice.

\subsubsection{Calibration Loss Weight ($\lambda_{\text{calib}}$)}

The calibration anchor weight $\lambda_{\text{calib}} = 10^{-4}$ was selected to balance two competing objectives: preventing calibration drift (which requires a non-zero anchor) while allowing the self-supervised flow loss to guide optimisation beyond AnyCalib's pseudo ground truth (which requires a sufficiently small anchor). The calibration accuracy results (Table~\ref{tab:calib-accuracy}) confirm that this value achieves this balance: our method improves upon AnyCalib on both Sintel (27.8\% vs 30.7\% MAPE) and TUM-RGBD (7.9\% vs 11.7\% MAPE), demonstrating that the anchor is loose enough to allow learning while tight enough to prevent divergence.

\subsubsection{KITTI Calibration Limitation}\label{sec:kitti-limitation}

While pose estimation improves on KITTI (Table~\ref{tab:pose-results}), calibration accuracy on KITTI degrades significantly compared to AnyCalib. KITTI uses automotive cameras with characteristics (long focal lengths, forward-facing motion, highway scenes) that differ substantially from the training data distribution. The MCT module, trained on indoor and urban walking datasets, does not generalise its calibration predictions to this out-of-distribution camera type. This suggests that extending the training data to include automotive scenarios would be necessary for robust KITTI calibration. Importantly, pose estimation still improves on KITTI despite the calibration limitation, indicating that the multi-frame consistency mechanism contributes independently of calibration quality.

\subsection{Qualitative Results}\label{sec:qualitative-results}

To visualise the effect of improved calibration and multi-frame consistency on trajectory estimation, we compare predicted camera trajectories against ground truth on a representative Sintel sequence (\texttt{market\_6}). This sequence was selected because it exhibits substantial camera translation through a dynamic market scene, making calibration quality directly visible in the trajectory shape.

Figure~\ref{fig:trajectory-comparison} shows the predicted trajectories after Sim(3) alignment. Our method (ATE $= 0.176$\,m) follows the ground truth path substantially more closely than AnyCam (ATE $= 0.352$\,m), reducing trajectory error by 2.0$\times$. The 3D view with camera frustums additionally shows that our method maintains more accurate camera orientations along the trajectory. The improvement is most visible in the translation direction: while AnyCam's trajectory drifts away from the ground truth path, our method maintains consistent alignment.

\begin{figure}[htpb]
    \centering
    \includegraphics[width=0.7\textwidth]{figures/trajectory_3d_side.png}
    \caption[Trajectory comparison]{3D trajectory comparison on Sintel \texttt{market\_6} after Sim(3) alignment, with camera frustums showing position and orientation accuracy. Our method (ATE $= 0.176$\,m) follows the ground truth more closely than AnyCam (ATE $= 0.352$\,m), a 2.0$\times$ improvement.}
    \label{fig:trajectory-comparison}
\end{figure}

Across all 23 Sintel sequences evaluated, our method achieves lower ATE than AnyCam on 13 sequences (57\%), with the largest improvements on sequences featuring substantial camera translation through dynamic scenes. This confirms that the combination of improved calibration and multi-frame consistency primarily benefits scenarios where translation accuracy is critical.

\section{Benchmark Implementation}\label{sec:benchmark-implementation}

A comprehensive benchmarking infrastructure was developed to enable systematic, reproducible evaluation.

\subsection{Automated Benchmark Scripts}

A suite of Python scripts implements the evaluation pipeline. Each script loads test sequences, runs inference, computes metrics (SO(3) error, translation direction, MAE/MAPE), and generates structured JSON results. Key scripts:

\begin{itemize}
    \item \texttt{benchmark\_phase\_c\_checkpoints.py}: Primary benchmark comparing our method against AnyCam (pose) and AnyCalib (calibration) baselines across Sintel, TUM-RGBD, and KITTI
    \item \texttt{generate\_thesis\_visualizations.py}: ATE-based trajectory scanning and 2D trajectory comparison generation
    \item \texttt{generate\_thesis\_3d\_trajectory.py}: 3D trajectory visualisation with camera frustums and optional point cloud reconstruction
    \item \texttt{generate\_thesis\_figures.py}: Per-frame focal length and error distribution plots
\end{itemize}

\subsection{Per-Epoch Monitoring}

During Phase~C training, a fixed validation set (50--100 samples) is evaluated every epoch for passive GT monitoring (not used in loss). Tracks training/validation loss components, rotation error, translation direction error, ATE, and calibration MAPE. Guides hyperparameter tuning and early stopping by detecting overfitting.

\section{Summary}

This chapter presented the evaluation framework and results:

\begin{enumerate}
    \item \textbf{Methodology}: Fully self-supervised training with GT used only for evaluation. Raw model comparisons without post-optimisation ensure fair assessment of the learned representations.

    \item \textbf{Metrics}: SO(3) rotation error, translation direction error, and ATE for pose evaluation; MAE/MAPE for calibration accuracy (when GT available); and focal length consistency CV as a pseudo-benchmark for datasets without GT.

    \item \textbf{Datasets}: MPI Sintel, TUM-RGBD, and KITTI provide ground truth for comprehensive evaluation across synthetic, indoor, and automotive scenarios.

    \item \textbf{Key findings}: The proposed method achieves improved pose accuracy compared to AnyCam (up to 31.7\% rotation improvement on TUM-RGBD) and more accurate calibration compared to AnyCalib (up to 32.5\% MAPE improvement on TUM-RGBD), confirming that feature-level multi-frame aggregation and staged self-supervised training effectively leverage the complementary strengths of both base systems. A KITTI calibration limitation was identified and discussed.

    \item \textbf{Design choices}: The optimal sequence length ($\text{max\_ahead}=3$) and calibration anchor weight ($\lambda_{\text{calib}} = 10^{-4}$) were validated in preliminary experiments and confirmed by full-scale evaluation results.
\end{enumerate}
